% Generated draft: requires mathematical and source review.
\documentclass[11pt]{article}
\usepackage[margin=25mm]{geometry}
\usepackage{fontspec}
\setmainfont{Latin Modern Roman}
\usepackage{amsmath,amssymb,mathtools,unicode-math}
\setmathfont{Latin Modern Math}
\usepackage{enumitem,needspace,xurl,hyperref}
\setlength{\parindent}{0pt}
\setlength{\parskip}{5pt}
\setlength{\emergencystretch}{3em}
\newcommand{\N}{\mathbb{N}}
\newcommand{\Z}{\mathbb{Z}}
\newcommand{\Q}{\mathbb{Q}}
\newcommand{\R}{\mathbb{R}}
\newcommand{\C}{\mathbb{C}}
\newcommand{\F}{\mathbb{F}}
\newcommand{\A}{\mathbb{A}}
\newcommand{\PP}{\mathbb{P}}
\newcommand{\E}{\mathbb{E}}
\newcommand{\eps}{\varepsilon}
\newcommand{\dd}{\,\mathrm{d}}
\begin{document}
\section*{UnsolvedMath ID: 1100207}
\textbf{Source problem code:} catalogue entry
\subsection{Definitions and mathematical statement}
For a group $G$, define the source's power-alternative property by
\[
\forall a,b\in G\ \exists n\in\Z_{>0}\quad
\bigl(\langle a^n,b^n\rangle\text{ is a free group}\bigr)
\ \lor\ 
\bigl(\langle a^n,b^n\rangle\text{ is abelian}\bigr).
\]
The same exponent $n$ is used for both elements and may depend on $G,a,b$. The literal source asks for a free subgroup, without explicitly requiring the displayed two generators to be a free basis of rank two. The abelian alternative means that all elements of that subgroup commute.

Wise asks whether every group $G$ acting properly discontinuously and cocompactly on a CAT(0) space has this property. In the usual geometric interpretation the action is by isometries, the space is geodesic and satisfies the CAT(0) comparison inequality, proper discontinuity requires that only finitely many translates of a compact set meet that set, and cocompactness means the orbit space is compact. The source does not explicitly specify its properness convention or any properness assumption on the metric space itself.

The parenthetical second question replaces this geometric hypothesis by the hypothesis that $G$ is automatic: in the usual sense, a finite generating set admits a regular language of representative words whose equality and multiplication-by-generator relations are recognized by finite automata reading padded word pairs synchronously. The question again asks for the displayed property for every pair $a,b\in G$.

The supplied research summary reports a negative answer to the CAT(0) question: Leary--Minasyan groups denoted $G_{k,m}$, with the stated restrictions
\[
-2m<k<2m,\qquad k\notin\{0,m,-m\},
\]
act properly and cocompactly on $\mathbb E^2\times T$ for a tree $T$ and contain $a,b$ with
\[
\forall n\in\Z_{>0},\qquad
\langle a^n,b^n\rangle\text{ is neither free nor abelian}.
\]
This is an attributed literature resolution: the record gives neither a presentation of these groups nor the construction of $a,b$ or the tree. It expressly says that this counterexample does not settle the automatic-group variant. Thus the overall solved label must not be read as a supplied solution of both questions.
\subsection{Short English statement}
Can taking a common positive power of any two elements always make their generated subgroup free or abelian, first in a CAT(0) group and separately in an automatic group? The supplied summary reports CAT(0) counterexamples of Leary and Minasyan. It explicitly leaves the automatic variant untouched by those counterexamples, so that variant is not resolved by the supplied solved label.
\subsection{Sources}
\begin{itemize}
\item Local source registry: \url{data/unsolvedmath/problems.json}, permanent UnsolvedMath ID $1100207$, source code catalogue entry, including the supplied research summary. Dataset attribution: Ulam AI and UnsolvedMath Contributors; original-source attribution and rights are retained.
\item Mladen Bestvina, \emph{Questions in Geometric Group Theory}, source-list date 2004, Question 2.7, attributed to D. Wise; the metadata identifies PDF page 7: \url{https://www.math.utah.edu/~bestvina/eprints/questions-updated.pdf}. The record reports manual transcription and access on July 29, 2026, and flags status and rights for review.
\item Leary--Minasyan (2021), Example 9.4 and Corollary 9.6 for the reported failure of the power alternative, and Corollary 9.3 for CAT(0), as cited by the research summary. No full title, group presentation, or separate link is supplied.
\item Martin (2024) and Hagen--Martin--Sartori (2025), cited in the summary as describing this as the first CAT(0) counterexample to the power alternative. Complete bibliographic details are not supplied.
\end{itemize}
\end{document}
